\documentclass[12pt,a4paper]{article}

\usepackage[utf8]{inputenc}
\usepackage[T1]{fontenc}
\usepackage[polish]{babel}
\usepackage{geometry}
\usepackage{array}
\usepackage{longtable}
\usepackage{hyperref}
\usepackage{enumitem}
\usepackage{graphicx}

\geometry{margin=2.5cm}
\setlength{\parindent}{0pt}
\setlength{\parskip}{6pt}

\title{Sprawozdanie z pracy projektowej\\\large Social Trip Planner / travel-app}
\author{Zespół projektowy}
\date{\today}

\begin{document}

\maketitle
\tableofcontents
\newpage

\section{Opis funkcjonalny systemu}

Projekt \textit{Social Trip Planner} jest aplikacją webową służącą do planowania podróży. System umożliwia użytkownikowi utworzenie konta, zalogowanie się, zarządzanie własnymi podróżami oraz korzystanie z publicznych inspiracji podróżniczych.

Aplikacja pozwala na:
\begin{itemize}
\item rejestrację i logowanie użytkownika,
\item uwierzytelnianie z użyciem tokenu JWT,
\item tworzenie własnych podróży,
\item wyświetlanie listy podróży użytkownika,
\item usuwanie podróży,
\item korzystanie z publicznych inspiracji podróżniczych,
\item kopiowanie inspiracji do własnych podróży,
\item obsługę listy uczestników podróży,
\item korzystanie z checklisty zadań,
\item przełączanie języka aplikacji między polskim i angielskim,
\item przełączanie motywu jasnego i ciemnego,
\item pobieranie dodatkowych informacji o kierunku podróży z zewnętrznych API.
\end{itemize}

System rozróżnia prywatne podróże użytkownika oraz publiczne inspiracje. Inspiracje są widoczne dla wszystkich użytkowników, natomiast prywatne podróże są przypisane do konkretnego konta. Dzięki temu użytkownik może przeglądać gotowe propozycje wyjazdów i skopiować wybraną inspirację do własnej listy podróży.

\section{Opis technologiczny}

Aplikacja została podzielona na część frontendową, backendową oraz bazę danych.

\subsection{Frontend}

Frontend został wykonany w technologii React z użyciem narzędzia Vite. Interfejs użytkownika korzysta z komponentów React, routingu po stronie klienta oraz asynchronicznej komunikacji z backendem przez Fetch API.

Wykorzystane technologie frontendowe:
\begin{itemize}
\item React,
\item Vite,
\item JavaScript,
\item SCSS,
\item React Router,
\item i18next.
\end{itemize}

React Router odpowiada za obsługę tras aplikacji, m.in. stron logowania, rejestracji i panelu użytkownika. Aplikacja korzysta również z mechanizmu chronionych tras, dzięki czemu panel użytkownika jest dostępny dopiero po zalogowaniu.

\subsection{Backend}

Backend został wykonany w języku Java 17 z użyciem frameworka Spring Boot. Aplikacja korzysta z architektury warstwowej: kontrolery obsługują żądania HTTP, serwisy realizują logikę biznesową, repozytoria odpowiadają za komunikację z bazą danych, a encje reprezentują dane systemu.

Wykorzystane technologie backendowe:
\begin{itemize}
\item Java 17,
\item Spring Boot,
\item Spring MVC,
\item Spring Security,
\item JWT,
\item Spring Data JPA,
\item Hibernate,
\item Maven,
\item Spring Cache,
\item SLF4J Logger.
\end{itemize}

Uwierzytelnianie zostało zrealizowane z użyciem tokenów JWT. Po poprawnym logowaniu użytkownik otrzymuje token, który jest następnie wysyłany w nagłówku żądań do zabezpieczonych endpointów.

\subsection{Baza danych}

Do przechowywania danych wykorzystano PostgreSQL. Schemat bazy danych jest zarządzany przy pomocy Flyway, co pozwala kontrolować kolejne migracje bazy danych i zachować spójność środowisk lokalnego oraz produkcyjnego.

W projekcie wykorzystano m.in. tabele:
\begin{itemize}
\item \texttt{users} -- użytkownicy systemu,
\item \texttt{trips} -- podróże,
\item \texttt{trip_user} -- relacja użytkowników i podróży,
\item \texttt{tasks} -- zadania/checklista,
\item \texttt{categories} -- kategorie podróży.
\end{itemize}

Mapowanie obiektowo-relacyjne zostało wykonane przy pomocy JPA/Hibernate.

\subsection{Integracje zewnętrzne}

Aplikacja korzysta z zewnętrznych API do pobierania dodatkowych informacji o kierunku podróży:
\begin{itemize}
\item zdjęcia kierunków podróży,
\item informacje pogodowe,
\item informacje walutowe.
\end{itemize}

Klucze API nie są przechowywane w repozytorium. Są przekazywane przez zmienne środowiskowe.

\subsection{Deployment}

Aplikacja została wdrożona w środowisku produkcyjnym:
\begin{itemize}
\item frontend: Vercel,
\item backend: Render,
\item baza danych: PostgreSQL na Renderze.
\end{itemize}

Adres wdrożonego frontendu:
\begin{center}
\url{https://travel-app-mocha-delta.vercel.app}
\end{center}

Adres backendu:
\begin{center}
\url{https://travel-app-api-lzua.onrender.com}
\end{center}

Wejście na główny adres backendu może zwracać komunikat \texttt{403 Forbidden}, ponieważ endpointy są zabezpieczone przez Spring Security. Właściwe endpointy API są dostępne pod ścieżką \texttt{/api}.

\section{Wdrożone zagadnienia kwalifikacyjne}

\begin{longtable}{|p{0.28\textwidth}|p{0.55\textwidth}|p{0.12\textwidth}|}
\hline
\textbf{Zagadnienie} & \textbf{Realizacja w projekcie} & \textbf{Status} \
\hline
Framework MVC & Backend wykonany w Spring Boot i Spring MVC. & Tak \
\hline
Framework CSS & Aplikacja korzysta ze stylowania SCSS oraz elementów Tailwind. & Tak \
\hline
Baza danych & Do projektu dołączono PostgreSQL. & Tak \
\hline
Cache & Zastosowano Spring Cache dla danych dotyczących kierunku podróży. & Tak \
\hline
Dependency manager & Backend korzysta z Maven, frontend z npm. & Tak \
\hline
HTML & Szkielet aplikacji generowany jest przez React i osadzony w HTML. & Tak \
\hline
CSS & Interfejs został ostylowany z użyciem SCSS. & Tak \
\hline
JavaScript & Frontend wykorzystuje JavaScript i React do obsługi interakcji. & Tak \
\hline
Routing & Zastosowano React Router oraz adresy typu pretty URLs. & Tak \
\hline
ORM & Wykorzystano JPA/Hibernate. & Tak \
\hline
Uwierzytelnianie & Zaimplementowano logowanie, rejestrację, hashowanie haseł oraz JWT. & Tak \
\hline
Lokalizacja & Aplikacja umożliwia przełączanie języka PL/EN z użyciem i18next. & Tak \
\hline
Mailing & Przygotowano serwis mailingowy po stronie backendu. Wysyłka wymaga poprawnej konfiguracji SMTP. & Częściowo \
\hline
Formularze & Zaimplementowano formularze logowania, rejestracji i dodawania podróży. & Tak \
\hline
Asynchroniczne interakcje & Frontend komunikuje się z backendem asynchronicznie przez Fetch API. & Tak \
\hline
Konsumpcja API & Aplikacja korzysta z zewnętrznych API dla zdjęć, pogody i walut. & Tak \
\hline
Publikacja API & Backend wystawia własne REST API. & Tak \
\hline
RWD & Interfejs jest responsywny i dostosowany do różnych szerokości ekranu. & Tak \
\hline
Logger & W backendzie użyto loggera do rejestrowania ważnych akcji systemowych. & Tak \
\hline
Deployment & Frontend wdrożono na Vercel, backend i bazę danych na Renderze. & Tak \
\hline
\end{longtable}

\section{Instrukcja lokalnego uruchomienia systemu}

\subsection{Wymagania}

Do lokalnego uruchomienia projektu wymagane są:
\begin{itemize}
\item Java 17,
\item Node.js,
\item npm,
\item Docker,
\item Git Bash lub inny terminal obsługujący skrypty Bash.
\end{itemize}

\subsection{Uruchomienie bazy danych}

W katalogu głównym projektu należy uruchomić:

\begin{verbatim}
docker compose up -d
\end{verbatim}

Nie należy używać polecenia:

\begin{verbatim}
docker compose down -v
\end{verbatim}

jeżeli lokalne dane w bazie mają zostać zachowane, ponieważ opcja \texttt{-v} usuwa wolumeny Dockera.

\subsection{Uruchomienie backendu}

Backend korzysta z pliku \texttt{backend/.env}, który nie jest commitowany do repozytorium. Plik powinien zawierać zmienne środowiskowe, m.in.:

\begin{verbatim}
SPRING_DATASOURCE_URL=jdbc:postgresql://localhost:5432/travel_app
SPRING_DATASOURCE_USERNAME=travel_user
SPRING_DATASOURCE_PASSWORD=your_password
JWT_SECRET=your_long_jwt_secret
PEXELS_API_KEY=your_pexels_key
WEATHER_API_KEY=your_weather_key
CURRENCY_API_KEY=your_currency_key
\end{verbatim}

Backend można uruchomić poleceniem:

\begin{verbatim}
./run-backend.sh
\end{verbatim}

Lokalny adres backendu:

\begin{verbatim}
http://localhost:8080
\end{verbatim}

\subsection{Uruchomienie frontendu}

W katalogu \texttt{frontend} należy wykonać:

\begin{verbatim}
npm install
npm run dev
\end{verbatim}

Lokalny adres frontendu:

\begin{verbatim}
http://localhost:5173
\end{verbatim}

Przykładowa zmienna środowiskowa dla frontendu:

\begin{verbatim}
VITE_API_URL=http://localhost:8080/api
\end{verbatim}

\section{Instrukcja zdalnego uruchomienia systemu}

Wdrożona aplikacja jest dostępna pod adresem:

\begin{center}
\url{https://travel-app-mocha-delta.vercel.app}
\end{center}

Backend działa jako usługa Render pod adresem:

\begin{center}
\url{https://travel-app-api-lzua.onrender.com}
\end{center}

W konfiguracji produkcyjnej frontend korzysta ze zmiennej:

\begin{verbatim}
VITE_API_URL=https://travel-app-api-lzua.onrender.com/api
\end{verbatim}

Backend na Renderze korzysta ze zmiennych środowiskowych:
\begin{itemize}
\item \texttt{SPRING_DATASOURCE_URL},
\item \texttt{SPRING_DATASOURCE_USERNAME},
\item \texttt{SPRING_DATASOURCE_PASSWORD},
\item \texttt{JWT_SECRET},
\item \texttt{PEXELS_API_KEY},
\item \texttt{WEATHER_API_KEY},
\item \texttt{CURRENCY_API_KEY},
\item \texttt{PORT=8080}.
\end{itemize}

Sekrety i klucze API nie są przechowywane w repozytorium.

\section{Przykładowe endpointy API}

\begin{longtable}{|p{0.18\textwidth}|p{0.35\textwidth}|p{0.37\textwidth}|}
\hline
\textbf{Metoda} & \textbf{Endpoint} & \textbf{Opis} \
\hline
POST & \texttt{/api/auth/register} & Rejestracja użytkownika. \
\hline
POST & \texttt{/api/auth/login} & Logowanie użytkownika. \
\hline
GET & \texttt{/api/trips} & Pobranie podróży użytkownika i inspiracji. \
\hline
POST & \texttt{/api/trips} & Dodanie nowej podróży. \
\hline
DELETE & \texttt{/api/trips/{id}} & Usunięcie podróży. \
\hline
POST & \texttt{/api/trips/{id}/copy} & Skopiowanie inspiracji do własnych podróży. \
\hline
GET & \texttt{/api/destinations/details?query=...} & Pobranie informacji o kierunku podróży. \
\hline
\end{longtable}

\section{Wnioski projektowe}

Projekt pozwolił na praktyczne połączenie wielu technologii webowych w jednej aplikacji. Najważniejszym elementem było zintegrowanie frontendu React z backendem Spring Boot oraz bazą danych PostgreSQL. Duże znaczenie miało także poprawne wdrożenie aplikacji na zewnętrznych platformach oraz konfiguracja zmiennych środowiskowych.

W trakcie pracy szczególnie istotne okazały się:
\begin{itemize}
\item poprawna konfiguracja CORS pomiędzy frontendem i backendem,
\item prawidłowe ustawienie adresu API w zmiennej \texttt{VITE_API_URL},
\item zabezpieczenie sekretów i kluczy API poza repozytorium,
\item zarządzanie schematem bazy danych przez Flyway,
\item rozróżnienie prywatnych podróży użytkownika i publicznych inspiracji.
\end{itemize}

Aplikacja spełnia główne założenia projektu i została wdrożona w środowisku produkcyjnym. Możliwe dalsze rozszerzenia obejmują m.in. rozbudowany panel administracyjny, pełną konfigurację wysyłki maili, testy automatyczne oraz dokumentację API w Swagger/OpenAPI.

\end{document}
